\section{Conventions}

\subsection{Metric}

There are two choices of metric, 
\begin{align}
    g_-=\textnormal{diag}\mqty(-&+&+&+),\ \ \ g_+=\textnormal{diag}\mqty(+&-&-&-)
\end{align}
I will always use $g_-$. This changes a lot of signs in a lot of places, as example,
\begin{align}
    \mathcal L = g_+^{\mu\nu}\frac12\partial_\mu \phi\partial_\nu\phi-V\qty[\phi]= -g_-^{\mu\nu}\frac12\partial_\mu \phi\partial_\nu\phi-V\qty[\phi]
\end{align}

Physics should be invariant under conventions, so they don't matter, as long as you know which conventions 
you are using and are consistent.

\subsection{Levi-Civita}

The Levi-Civita is the (unique) totally antisymmetric four-pseudo-tensor, $\epsilon_{\alpha\beta\mu\nu}$, there are two possible 
choices of normalization,
\begin{align}
    \epsilon^+_{0123}=1,\ \ \ \epsilon^-_{0123}=-1
\end{align}
I will always use $\epsilon^+$. We define also the spatial Levi-Civita by restricting the covariant one,
\begin{align}
    \epsilon_{ijk} \coloneq \epsilon_{0ijk}
\end{align}